% META: {"id": "TCOMPL-ME-CR-013", "chapitre": "TCOMPL-MODELES-EVOLUTION", "type_objet": "cours", "capacites_codes": ["C6"], "sources_inspiration": [], "mode_creation": "ex_nihilo", "status": "needs_review"}

\section{Limite d'une suite : notion, opérations, comparaison}

% ============================================================
% STRATE 1 — Notion intuitive
% ============================================================

\definition[1]{
Soit $(u_n)$ une suite de nombres réels.

$\bullet$ $(u_n)$ a pour \textbf{limite finie} $\ell$ lorsque les termes $u_n$ finissent par s'accumuler autour de $\ell$, aussi pres qu'on le souhaite, des que $n$ est assez grand. On note $\displaystyle\lim_{n \to +\infty} u_n = \ell$ et l'on dit que la suite \textbf{converge} vers $\ell$.

$\bullet$ $(u_n)$ a pour \textbf{limite} $+\infty$ lorsque $u_n$ finit par depasser n'importe quel nombre fixe a l'avance, des que $n$ est assez grand. On note $\displaystyle\lim_{n \to +\infty} u_n = +\infty$, et l'on dit que la suite \textbf{diverge} vers $+\infty$ (définition analogue pour $-\infty$).
}

\margeAppui{
Une suite qui ne converge pas ne tend pas forcement vers l'infini : la suite $u_n = (-1)^n$ oscille entre $-1$ et $1$ sans jamais s'accumuler autour d'une seule valeur, et sans depasser aucun seuil. Elle n'a pas de limite.
}

\exemple{
La suite recurrente $u_{n+1} = 0{,}5\,u_n + 3$ avec $u_0=0$, resolue plus haut, a pour terme general $u_n = 6 - 6\times0{,}5^n$. Comme $0{,}5^n$ tend vers $0$, la suite converge vers sa valeur d'equilibre $6$ : le modèle se stabilise.
}

% ============================================================
% STRATE 2 — Operations sur les limites
% ============================================================

\propriete[Opérations sur les limites]{
Soit $(u_n)$ et $(v_n)$ deux suites de limites respectives $\ell$ et $\ell'$ (finies ou infinies). Lorsque le résultat n'est pas une forme indeterminee :
\begin{itemize}
  \item \textbf{Somme} : $u_n+v_n$ tend vers $\ell+\ell'$. Cas infini : $\ell$ fini et $\ell'=+\infty$ donnent $+\infty$.
  \item \textbf{Produit} : $u_n\,v_n$ tend vers $\ell\times\ell'$. Cas infini : $\ell>0$ et $\ell'=+\infty$ donnent $+\infty$.
  \item \textbf{Quotient} : si $\ell' \neq 0$, $\dfrac{u_n}{v_n}$ tend vers $\dfrac{\ell}{\ell'}$ ; si $\ell$ est fini et $\ell' = \pm\infty$, le quotient tend vers $0$.
\end{itemize}
}

\margeAppui{
Quatre situations echappent a ces règles et s'appellent \textbf{formes indeterminees} : $(+\infty)+(-\infty)$, $0\times\infty$, $\dfrac{\infty}{\infty}$ et $\dfrac{0}{0}$. Il faut alors transformer l'ecriture — factoriser par le terme dominant, par exemple — avant de conclure.
}

\exemple{
$u_n = \dfrac{3n+5}{n+2}$ presente la forme indeterminee $\dfrac{\infty}{\infty}$. En factorisant par $n$ au numérateur et au dénominateur :
\[
  u_n = \frac{n\left(3+\frac{5}{n}\right)}{n\left(1+\frac{2}{n}\right)} = \frac{3+\frac{5}{n}}{1+\frac{2}{n}} \xrightarrow[n \to +\infty]{} \frac{3+0}{1+0} = 3 .
\]
}

% ============================================================
% STRATE 3 — Comparaison et gendarmes
% ============================================================

\propriete[Passage a la limite dans les inégalités]{
Si $u_n \leq v_n$ pour tout $n$ assez grand, et si $(u_n)$ et $(v_n)$ convergent vers $\ell$ et $\ell'$, alors $\ell \leq \ell'$.
}

\margeAppui{
Attention : une inégalité \textbf{stricte} entre les termes ne se transmet pas a la limite. On a $\dfrac{1}{n+1} < \dfrac{1}{n}$ pour tout $n \geq 1$, et pourtant les deux suites ont la \emph{même} limite $0$ : le passage a la limite ne conserve que l'inégalité large.
}

\theoreme[Théorème des gendarmes]{
Si $u_n \leq w_n \leq v_n$ pour tout $n$ assez grand, et si $(u_n)$ et $(v_n)$ convergent vers la \textbf{même} limite $\ell$, alors $(w_n)$ converge aussi vers $\ell$.
}

\exemple{
Soit $w_n = \dfrac{\sin n}{n}$ pour $n \geq 1$. Comme $-1 \leq \sin n \leq 1$, on encadre :
\[
  -\frac{1}{n} \leq w_n \leq \frac{1}{n}.
\]
Les deux suites encadrantes tendent vers $0$ : d'après le théorème des gendarmes, $w_n$ tend vers $0$.
}

% BEGIN-VERIFY
% from sympy import symbols, limit, oo, sin, Rational, sqrt, simplify
% n = symbols('n', positive=True, integer=True)
% # Le modele d'evolution du chapitre tend bien vers sa valeur d'equilibre.
% u = 6 - 6 * Rational(5, 10)**n
% assert limit(u, n, oo) == 6
% # Le terme general reproduit bien la recurrence u_{n+1} = 0.5 u_n + 3.
% terme = lambda k: 6 - 6 * Rational(1, 2)**k
% for k in range(8):
%     assert terme(k + 1) == Rational(1, 2) * terme(k) + 3
% assert terme(0) == 0
% # La forme indeterminee infini sur infini, apres factorisation.
% assert limit((3*n + 5) / (n + 2), n, oo) == 3
% # L'encadrement des gendarmes est vrai terme a terme sur un echantillon.
% import math
% for k in range(1, 200):
%     assert -1/k <= math.sin(k)/k <= 1/k
% # Les deux gendarmes ont la meme limite, donc la suite encadree aussi.
% assert limit(-1/n, n, oo) == 0 and limit(1/n, n, oo) == 0
% assert limit(sin(n)/n, n, oo) == 0
% # Une inegalite stricte ne se transmet pas : memes limites malgre 1/(n+1) < 1/n.
% for k in range(1, 200):
%     assert 1/(k+1) < 1/k
% assert limit(1/(n+1), n, oo) == limit(1/n, n, oo)
% # (-1)^n n'a pas de limite : ses deux sous-suites different.
% assert (-1)**200 == 1 and (-1)**201 == -1
% END-VERIFY

\erreurFrequente{
\textbf{Appliquer le théorème des gendarmes avec deux limites différentes.} L'encadrement $-\dfrac{2}{n} \leq w_n \leq \dfrac{1}{n}$ suffit a conclure car les deux bornes tendent vers $0$ ; mais un encadrement $1-\dfrac{1}{n} \leq w_n \leq 2+\dfrac{1}{n}$ ne permet rien de conclure : les gendarmes tendent vers $1$ et vers $2$, et la suite encadree peut osciller entre les deux.
}
